\section{Contributions}

This dissertation, \emph{Defending Emerging Software Systems in Untrusted Environments}, 
studies a common security problem across two families of 
increasingly autonomous software systems. 
Smart contracts already execute valuable financial logic with 
no human intervention and must tolerate arbitrary 
on-chain interaction. 
LLM-based coding agents increasingly translate vague specifications 
in natural language
into concrete code while relying on opaque models and weakly 
trusted online information. 
In both cases, harmful behavior may be induced by factors that cannot 
be fully enumerated or validated in advance. This thesis 
therefore studies how to constrain harmful behavior 
of a software system under untrusted environments
such as 
arbitrary user inputs, 
adversarial interactions, opaque model behavior, 
and weakly trusted online information.


We organize the technical contribution as four connected layers. 
Layers I and II focus on smart contracts, where 
the main challenge is runtime protection 
under arbitrary adversarial interaction. 
Layers III and IV focus on LLM-centered coding systems, 
where the main challenge is that preventing malicious code generation
before it gets executed.  
Together, these layers form one defense story: 
how to secure high-stake software systems that 
operate in untrusted environments. 


This thesis makes the following contributions.
\begin{itemize}
    \item \textbf{A unified defense framing for emerging software 
    systems in untrusted environments.}
    We formulate a common perspective for software systems whose behavior 
    may be shaped by arbitrary user actions,
    adversarial interactions, opaque models, contaminated training data, 
    or weakly trusted online information.

    % \item \textbf{Layer I: Runtime invariants for smart contracts.}
    % We study how to learn and enforce runtime invariants, namely rules 
    % that should continue to hold during benign executions,
    % in order to stop abnormal transaction behaviors in deployed financial software.

    % \item \textbf{Layer II: Control-flow integrity for smart contracts.}
    % We study how to constrain unexpected invocation patterns across 
    % functions and contracts,
    % thereby protecting deployed smart contracts against attacks that exploit 
    % unintended interaction structure.

    % \item \textbf{Layer III: Auditing harmful artifacts in production LLMs.}
    % We study to what extent production models can emit harmful artifacts in ordinary 
    % developer workflows,
    % providing an empirical view of upstream risk before generated 
    % code is trusted or deployed.

    % \item \textbf{Layer IV: Search- and documentation-level poisoning of coding agents.}
    % We study how weakly trusted external content can poison 
    % coding agents through retrieval and reasoning,
    % and how practical defenses can reduce this risk before 
    % harmful actions are taken.

    \item \noindent \textbf{Layer I: Demystifying invariant effectiveness for securing smart contracts.}
    it proposes runtime invariant generation  and enforcement techniques
    for smart contract security, where invariants are rules that 
    should continue to hold during normal execution.

    \item \noindent \textbf{Layer II: Enforcing control flow integrity on DeFi smart contracts.}
    it proposes algorithms for control flow integrity 
    enforcement across contract interactions, 
    with minimal runtime overhead; informally, 
    this layer ensures that contract functions are called 
    in legitimate patterns.

    \item \noindent \textbf{Layer III: Auditing malicious scam endpoints in production LLMs.}
    it audits LLMs when generating code under realistic developer-style prompts, 
    and finds that LLMs can generate malicious scam URLs with high frequency. 
    This layer reveals a new attack vector in LLMs, and proposes 
    a practical detection system to audit malicious endpoints in the pre-training 
    data of LLMs.

    \item \noindent \textbf{Layer IV: Defending Search- and Documentation-Level Poisoning of Coding Agents.}
    it studies how poisoned external content retrieved by coding agents 
    can propagate into harmful coding outcomes, and how to defend against this pipeline.




\end{itemize}
